\documentclass[11pt]{article}
\usepackage[margin=1in]{geometry}
\usepackage{lmodern}
\usepackage{enumitem}
\usepackage{hyperref}
\setlength{\parskip}{0.6em}
\setlength{\parindent}{0pt}
\title{Machine Learning --- Contribution Information Sheet}
\author{}
\date{}
\begin{document}
\maketitle

\textbf{1. What is the main claim of the paper? Why is this an important contribution to the machine learning literature?}

\textbf{Main claim.} In a preregistered experiment, iterated self-revision of a written constitution did not selectively erode language models' commitments to human oversight. Erosion concentrated instead in stylistic self-restraints and was driven by reflection rather than by repeated editing. This runs against a long-standing argument that corrigibility is not reflectively stable.

\textbf{Why it matters for machine learning.} Models are increasingly trained and deployed with procedures that let them critique and revise their own principles: constitutional training, self-refinement and self-evolving agents. To our knowledge, this is the first controlled, preregistered measurement of what such revision does to different kinds of stated commitments. It also provides a reusable method for making such measurements: iterated transmission used as a measurement instrument, matched item categories with control conditions, discrete-time survival analysis, and a two-judge reliability check.

\textbf{2. What is the evidence you provide to support your claim?}

\textbf{Design.} 7 open-weight model configurations $\times$ 4 conditions $\times$ 25 chains $\times$ 20 rounds (700 chains; 431,359 at-risk item-rounds; 2,342 erosion events), with paraphrase and neutral-edit control conditions and matched oversight and agentic items. Hypotheses, sample size and analysis code were preregistered and frozen, and the code was validated on simulated data before data collection. A second, independent LLM judge double-coded a stratified 25\% sample (Krippendorff's $\alpha = 0.733$).

\textbf{Results.} The erosion hazard ratio for oversight vs.\ matched agentic items was 0.495 (95\% CI 0.365--0.671). We report its fragility openly: models that include between-item variation give intervals that include 1, and the result depends on particular items. Stylistic items eroded at 43.5\% under self-reflection, against 7.8\% under neutral editing. Revising one's own constitution did not differ from revising another assistant's. Installing self-revised oversight clauses had no behavioural effect (COR-swap minus AGENT-swap 0.26, 95\% CI $-$0.21 to 0.73), in a battery where positive controls lowered acceptance of legitimate oversight by up to 0.95.

\textbf{3. What papers by other authors make the most closely related contributions, and how is your paper related to them?}
\begin{itemize}[leftmargin=*]
\item \textbf{Soares et al.\ (2015)} state the theoretical argument that corrigibility is not reflectively stable; we test it empirically.
\item \textbf{Bai et al.\ (2022)} train models against human-written principles (Constitutional AI); we let models rewrite their own principles and measure what survives.
\item \textbf{S.\ Das et al.\ (2026)} and \textbf{Han et al.\ (2025)} show that safety constraints erode in self-evolving agents under task or reward pressure; we remove pressure, reward and memory to isolate unpressured reflection, and we test selectivity against matched categories.
\item \textbf{Greenblatt et al.\ (2024), Meinke et al.\ (2024)} and \textbf{Schlatter et al.\ (2025)} measure oversight-relevant behaviour under goal conflict; we study stated commitments under reflection and test whether changes to them carry over into behaviour.
\item \textbf{J.\ Perez et al.\ (2024)} and \textbf{Wu et al.\ (2026)} use iterated transmission and self-refinement to reveal model attractors; we use this effect as a measurement instrument for values.
\item Two concurrent exploratory projects (\textbf{heyuday, 2026; Wu, 2026}) iterate constitution editing but include no control categories, no category-level analysis and no behavioural follow-up.
\end{itemize}

\textbf{4. Have you published parts of this work elsewhere? If so, how does this paper differ?}

No part of this work has been published, posted as a preprint, or submitted elsewhere.

\end{document}
